% Options for packages loaded elsewhere
\PassOptionsToPackage{unicode}{hyperref}
\PassOptionsToPackage{hyphens}{url}
\PassOptionsToPackage{dvipsnames,svgnames,x11names}{xcolor}
\documentclass[
]{article}
\usepackage{xcolor}
\usepackage[margin = 0.5in]{geometry}
\usepackage{amsmath,amssymb}
\setcounter{secnumdepth}{-\maxdimen} % remove section numbering
\usepackage{iftex}
\ifPDFTeX
  \usepackage[T1]{fontenc}
  \usepackage[utf8]{inputenc}
  \usepackage{textcomp} % provide euro and other symbols
\else % if luatex or xetex
  \usepackage{unicode-math} % this also loads fontspec
  \defaultfontfeatures{Scale=MatchLowercase}
  \defaultfontfeatures[\rmfamily]{Ligatures=TeX,Scale=1}
\fi
\usepackage{lmodern}
\ifPDFTeX\else
  % xetex/luatex font selection
\fi
% Use upquote if available, for straight quotes in verbatim environments
\IfFileExists{upquote.sty}{\usepackage{upquote}}{}
\IfFileExists{microtype.sty}{% use microtype if available
  \usepackage[]{microtype}
  \UseMicrotypeSet[protrusion]{basicmath} % disable protrusion for tt fonts
}{}
\makeatletter
\@ifundefined{KOMAClassName}{% if non-KOMA class
  \IfFileExists{parskip.sty}{%
    \usepackage{parskip}
  }{% else
    \setlength{\parindent}{0pt}
    \setlength{\parskip}{6pt plus 2pt minus 1pt}}
}{% if KOMA class
  \KOMAoptions{parskip=half}}
\makeatother
\usepackage{graphicx}
\makeatletter
\newsavebox\pandoc@box
\newcommand*\pandocbounded[1]{% scales image to fit in text height/width
  \sbox\pandoc@box{#1}%
  \Gscale@div\@tempa{\textheight}{\dimexpr\ht\pandoc@box+\dp\pandoc@box\relax}%
  \Gscale@div\@tempb{\linewidth}{\wd\pandoc@box}%
  \ifdim\@tempb\p@<\@tempa\p@\let\@tempa\@tempb\fi% select the smaller of both
  \ifdim\@tempa\p@<\p@\scalebox{\@tempa}{\usebox\pandoc@box}%
  \else\usebox{\pandoc@box}%
  \fi%
}
% Set default figure placement to htbp
\def\fps@figure{htbp}
\makeatother
\setlength{\emergencystretch}{3em} % prevent overfull lines
\providecommand{\tightlist}{%
  \setlength{\itemsep}{0pt}\setlength{\parskip}{0pt}}
\usepackage{color}
\usepackage{framed}
\setlength{\fboxsep}{.8em}

\usepackage{tcolorbox}

\tcbset{
  mystyle/.style={
    before skip=-10pt, % Adjust this value
    after skip=-10pt   % Adjust this value
  }
}
\newtcolorbox{mynoteblock}[1][]{mystyle, #1}

\newtcolorbox{blackbox}{
  colback=black,
  colframe=orange,
  coltext=white,
  boxsep=5pt,
  arc=4pt}
  
\newtcolorbox{ltbluebox}{
  colback=blue!5!white,
  colframe=blue!75!black,
  boxsep=5pt,
  arc=4pt}
  
\newtcolorbox{ltgreenbox}{
  colback=green!5!white,
  colframe=blue!75!black,
  boxsep=5pt,
  arc=4pt}
  
\newtcolorbox{ltpurplebox}{
  colback=purple!5!white,
  colframe=blue!75!black,
  boxsep=5pt,
  arc=4pt}
  


\newenvironment{cols}[1][]{}{}

\newenvironment{col}[1]{\begin{minipage}[t]{#1}\ignorespaces\vspace{0pt}}{%
\end{minipage}
\ifhmode\unskip\fi
\aftergroup\useignorespacesandallpars}

\def\useignorespacesandallpars#1\ignorespaces\fi{%
#1\fi\ignorespacesandallpars}

\makeatletter
\def\ignorespacesandallpars{%
  \@ifnextchar\par
    {\expandafter\ignorespacesandallpars\@gobble}%
    {}%
}
\makeatother

\usepackage{awesomebox}
\setlength{\aweboxvskip}{1mm}


\pagenumbering{gobble}

\usepackage{enumitem}
\setlist{itemsep=1pt, parsep=0pt} %# Adjust vertical spacing
\setlist[itemize,1]{leftmargin=*, labelsep=0.3em} %# Adjust horizontal spacing
\setlist[enumerate,1]{leftmargin=*, labelsep=0.3em} %# Adjust horizontal spacing

\usepackage{multirow}
\usepackage{multicol}
\usepackage{colortbl}
\usepackage{hhline}
\newlength\Oldarrayrulewidth
\newlength\Oldtabcolsep
\usepackage{longtable}
\usepackage{array}
\usepackage{hyperref}
\usepackage{float}
\usepackage{wrapfig}
\usepackage{bookmark}
\IfFileExists{xurl.sty}{\usepackage{xurl}}{} % add URL line breaks if available
\urlstyle{same}
\hypersetup{
  colorlinks=true,
  linkcolor={Maroon},
  filecolor={Maroon},
  citecolor={Blue},
  urlcolor={blue},
  pdfcreator={LaTeX via pandoc}}

\author{}
\date{\vspace{-2.5em}}

\begin{document}

\vspace{-1cm}

\section{Catch/Quota Setting Step 4: ABC
Setting}\label{catchquota-setting-step-4-abc-setting}

\begin{noteblock}
Catch and Quota Setting is a multistep process involving decisions by
the Northeast Regional Coordinating Committee (NRCC; including the Mid
Atlantic and New England Councils, Atlantic States Marine Fisheries
Commission, and NOAA's Greater Atlantic Regional Fisheries Office and
Northeast Fisheries Science Center), assessment scientists, assessment
peer reviewers, the Council's Scientific and Statistical Committee
(SSC), Monitoring Committees, and the Council. Information for each
stock and fishery is reviewed and synthesized to determine catch and
quota levels that meet management objectives.

\end{noteblock}

\begin{itemize}
\tightlist
\item
  \textbf{Key question: What additional indicators could inform SSC
  decisions at the ABC Determination step in this process? Could
  habitat, ecosystem, or socioeconomic indicators improve the process?}
\item
  \textbf{Scope: Should we develop a general process for using
  indicators in the ABC Determination step of the catch/quota setting
  process, or focus on a specific application for a single species?
  Which species?}
\end{itemize}

\begin{cols}

\begin{col}{0.49\textwidth}

\begin{ltbluebox}

\subsection{What information is used
now?}\label{what-information-is-used-now}

\begin{ltpurplebox}
\textbf{4. Determine Acceptable Biological Catch (ABC)}

The SSC uses reviewed assessment outputs within a structured framework
to determine ABC from assessment uncertainty based on input data, model
diagnostics and performance, and ecosystem considerations.

\end{ltpurplebox}

The MT assessment gives current status and short term projections of
overfishing limit (OFL) that may or may not already include ecosystem
information. The SSC may discuss performance of short term projections
or terminal year estimates in the assessment; ecosystem information can
inform decisions about projections or terminal year estimates. The SSC
may request new OFL estimates/projections from the assessment lead based
on this discussion.

\vspace{1ex}

The SSC uses its
\href{https://www.mafmc.org/s/Final-OFL-CV-guidance-document_06_24.pdf}{OFL
CV process} to evaluate uncertainty arising from data quality, model
selection/identification, retrospective analysis, empirical comparisons,
ecosystem factors, and recruitment stanzas. Ecosystem information can
feed into:

\begin{itemize}
\tightlist
\item
  Ecosystem factors:
  \href{https://journals.plos.org/plosone/article?id=10.1371/journal.pone.0146756}{Climate
  Vulnerability Analysis} results are currently included for each stock,
  can include multiple other environmentally driven changes in the stock
\item
  Recruitment stanzas: links to environmental drivers can decrease
  recruitment uncertainty
\end{itemize}

Ecosystem information that is presented along with the assessment is
easiest to use in the OFL CV determination, but nothing limits the SSC
from using ecosystem information presented from another source. The SSC
has based decisions on published literature in the past.

\vspace{1ex}

The OFL CV is used along with the estimate of current stock status in
the
\href{http://www.mafmc.org/s/MAFMC-ABC-Control-Rule-White-Paper.pdf}{MAFMC
ABC control rule} to determine the ABC that is recommended to the
Council. The SSC often provides annually averaged constant ABC for a
specified period (2-3 years) as well as annual ABC for each year at the
request of the Council.

\end{ltbluebox}

\end{col}

\begin{col}{0.02\textwidth}
~

\end{col}

\begin{col}{0.49\textwidth}

\begin{ltgreenbox}

\subsection{Information sources for
indicators}\label{information-sources-for-indicators}

\href{https://www.fisheries.noaa.gov/new-england-mid-atlantic/ecosystems/state-ecosystem-reports-northeast-us-shelf}{State
of the Ecosystem},
\href{https://www.mafmc.org/s/2_MAB_RiskAssess_2024.pdf}{Risk
Assessment}, and
\href{https://www.fisheries.noaa.gov/new-england-mid-atlantic/ecosystems/ecosystem-and-socioeconomic-profile-development-and-reports}{Ecosystem
and Socioeconomic Profiles} conceptual models and indicators for MAFMC
stocks could inform SSC discussion of terminal year estimates, short
term projection performance, and assessment uncertainty.

\vspace{1ex}

Similar indicators as above plus habitat indicators
\href{https://www.fisheries.noaa.gov/resource/map/essential-fish-habitat-mapper}{EFH
Habitat Mapper}, \href{https://nrha.shinyapps.io/dataexplorer/}{NRHA
Data Explorer} could further inform OFL CV discussions. Systematic
translation of stock specific indicators into risk criteria as in the
Pacific could simplify OFL CV discussions.

\global\setlength{\Oldarrayrulewidth}{\arrayrulewidth}

\global\setlength{\Oldtabcolsep}{\tabcolsep}

\setlength{\tabcolsep}{2pt}

\renewcommand*{\arraystretch}{1.5}



\providecommand{\ascline}[3]{\noalign{\global\arrayrulewidth #1}\arrayrulecolor[HTML]{#2}\cline{#3}}

\begin{longtable}[c]{|p{0.60in}|p{2.50in}}

\caption{Pacific\ Pilot\ Risk\ Levels\ and\ Indicator\ Criteria}\\

\hhline{>{\arrayrulecolor[HTML]{666666}\global\arrayrulewidth=1.5pt}->{\arrayrulecolor[HTML]{666666}\global\arrayrulewidth=1.5pt}-}

\multicolumn{1}{>{\raggedright}m{\dimexpr 0.6in+0\tabcolsep}}{\textcolor[HTML]{000000}{\fontsize{9}{9}\selectfont{Risk\ Level}}} & \multicolumn{1}{>{\raggedright}m{\dimexpr 2.5in+0\tabcolsep}}{\textcolor[HTML]{000000}{\fontsize{9}{9}\selectfont{Ecosystem\ and\ Environmental\ Conditions}}} \\

\noalign{\global\arrayrulewidth 0pt}\arrayrulecolor[HTML]{000000}

\hhline{>{\arrayrulecolor[HTML]{666666}\global\arrayrulewidth=1.5pt}->{\arrayrulecolor[HTML]{666666}\global\arrayrulewidth=1.5pt}-}\endfirsthead \caption[]{Pacific\ Pilot\ Risk\ Levels\ and\ Indicator\ Criteria}\\

\hhline{>{\arrayrulecolor[HTML]{666666}\global\arrayrulewidth=1.5pt}->{\arrayrulecolor[HTML]{666666}\global\arrayrulewidth=1.5pt}-}

\multicolumn{1}{>{\raggedright}m{\dimexpr 0.6in+0\tabcolsep}}{\textcolor[HTML]{000000}{\fontsize{9}{9}\selectfont{Risk\ Level}}} & \multicolumn{1}{>{\raggedright}m{\dimexpr 2.5in+0\tabcolsep}}{\textcolor[HTML]{000000}{\fontsize{9}{9}\selectfont{Ecosystem\ and\ Environmental\ Conditions}}} \\

\noalign{\global\arrayrulewidth 0pt}\arrayrulecolor[HTML]{000000}

\hhline{>{\arrayrulecolor[HTML]{666666}\global\arrayrulewidth=1.5pt}->{\arrayrulecolor[HTML]{666666}\global\arrayrulewidth=1.5pt}-}\endhead



\multicolumn{1}{>{\cellcolor[HTML]{00FF00}\raggedright}m{\dimexpr 0.6in+0\tabcolsep}}{\textcolor[HTML]{000000}{\fontsize{9}{9}\selectfont{Level\ 1:\ Favorable}}} & \multicolumn{1}{>{\cellcolor[HTML]{00FF00}\raggedright}m{\dimexpr 2.5in+0\tabcolsep}}{\textcolor[HTML]{000000}{\fontsize{9}{9}\selectfont{Indicators\ not\ used\ in\ the\ stock\ assessment\ show\ medium\ to\ high\ level\ of\ agreement\ and\ moderate\ to\ strong\ evidence\ supporting\ high\ species\ productivity.}}} \\

\noalign{\global\arrayrulewidth 0pt}\arrayrulecolor[HTML]{000000}





\multicolumn{1}{>{\cellcolor[HTML]{FFFFFF}\raggedright}m{\dimexpr 0.6in+0\tabcolsep}}{\textcolor[HTML]{000000}{\fontsize{9}{9}\selectfont{Level\ 2:\ Neutral}}} & \multicolumn{1}{>{\cellcolor[HTML]{FFFFFF}\raggedright}m{\dimexpr 2.5in+0\tabcolsep}}{\textcolor[HTML]{000000}{\fontsize{9}{9}\selectfont{Majority\ of\ indicators\ show\ no\ notable\ trends\ and/or\ no\ apparent\ environmental\ and\ ecosystem\ concerns.}}} \\

\noalign{\global\arrayrulewidth 0pt}\arrayrulecolor[HTML]{000000}





\multicolumn{1}{>{\cellcolor[HTML]{FF0000}\raggedright}m{\dimexpr 0.6in+0\tabcolsep}}{\textcolor[HTML]{000000}{\fontsize{9}{9}\selectfont{Level\ 3:\ Unfavorable}}} & \multicolumn{1}{>{\cellcolor[HTML]{FF0000}\raggedright}m{\dimexpr 2.5in+0\tabcolsep}}{\textcolor[HTML]{000000}{\fontsize{9}{9}\selectfont{Majority\ of\ indicators\ show\ medium\ to\ high\ level\ of\ agreement\ and\ moderate\ to\ strong\ evidence\ supporting\ low\ species\ productivity}}} \\

\noalign{\global\arrayrulewidth 0pt}\arrayrulecolor[HTML]{000000}

\hhline{>{\arrayrulecolor[HTML]{666666}\global\arrayrulewidth=1.5pt}->{\arrayrulecolor[HTML]{666666}\global\arrayrulewidth=1.5pt}-}



\end{longtable}



\arrayrulecolor[HTML]{000000}

\global\setlength{\arrayrulewidth}{\Oldarrayrulewidth}

\global\setlength{\tabcolsep}{\Oldtabcolsep}

\renewcommand*{\arraystretch}{1}

\subsubsection{Potential Pathways}\label{potential-pathways}

\begin{itemize}
\tightlist
\item
  Refine OFL CV criteria for considering integrated ecosystem conditions
\item
  Develop criteria for considering ecosystem impacts in short term
  projections
\end{itemize}

\begin{ltpurplebox}

\subsection{Information gaps}\label{information-gaps}

Ecosystem and Socioeconomic Profiles are currently unavailable for
summer flounder, scup, mackerel, butterfish, surfclam, ocean quahog,
blueline tilefish, spiny dogfish, and monkfish.

\end{ltpurplebox}

\end{ltgreenbox}

\end{col}

\end{cols}

\end{document}
